\chapter{Recerca prèvia}
\label{c:recerca previa}

Abans de començar a programar i a fer proves amb dades reals, vaig dedicar temps a informar-me sobre el context del meu treball. Aquesta recerca prèvia m’ha servit per entendre millor el funcionament dels serveis de bicicleta compartida i el paper que poden tenir dins la mobilitat sostenible.

\section{El Bicing i la mobilitat urbana}
El \textbf{Bicing} és el sistema de bicicletes compartides de Barcelona, creat l’any 2007. Des de llavors s’ha convertit en un element habitual de la ciutat: milers de persones l’utilitzen cada dia per desplaçar-se de manera ràpida i sense contaminació. El seu objectiu és reduir l’ús del cotxe i oferir una alternativa més sana i respectuosa amb el medi ambient. Avui en dia compta amb centenars d’estacions i bicicletes, tant mecàniques com elèctriques.

Un aspecte destacable és que el servei és digital: les estacions informen en temps real de quantes bicicletes hi ha disponibles i quants espais lliures queden. Això permet que els usuaris puguin consultar l’estat des del mòbil abans d’anar-hi.

\section{El concepte de dades obertes}
Moltes ciutats, inclosa Barcelona, comparteixen informació dels seus serveis amb la ciutadania. Aquest concepte rep el nom de \textit{dades obertes} (\textit{open data}) i consisteix a posar a l’abast de tothom informació que pot ser útil per fer estudis, aplicacions o projectes de recerca.

En el cas del Bicing, aquestes dades es poden consultar a través d’un mecanisme senzill i públic: una API oberta. Gràcies a això, qualsevol persona amb coneixements de programació pot descarregar-se l’estat de les estacions i treballar-hi.

\section{Què és una API?}
Una \textbf{API}~\cite{api-def} (sigles d’\textit{Application Programming Interface}) és, en paraules simples, com una porta d’entrada que permet a dos programes parlar entre ells. En lloc de mirar directament dins la base de dades del Bicing, els programadors demanen la informació a través d’aquesta “porta”. L’API respon amb un paquet de dades que conté la informació actualitzada de les estacions.

Això fa que sigui fàcil i segur utilitzar aquestes dades sense necessitat de conèixer els detalls tècnics interns del sistema.

\section{El format JSON}
Les dades que s’obtenen de l’API del Bicing es presenten en \textbf{format JSON}~\cite{json-org}, unes sigles que volen dir \textit{JavaScript Object Notation}. Aquest format és com una llista de parelles \emph{clau–valor}. Per exemple, una estació pot aparèixer així:

\begin{verbatim}
{
  "name": "Gran Via - Marina",
  "free_bikes": 5,
  "empty_slots": 7
}
\end{verbatim}

En aquest cas veiem que l’estació “Gran Via - Marina” té 5 bicicletes disponibles i 7 espais lliures. Aquest format és molt comú perquè és fàcil de llegir i d’entendre, tant per les persones com pels ordinadors.

\section{El llenguatge Python}
Per desenvolupar el meu projecte vaig escollir \textbf{Python}~\cite{python-docs}, un llenguatge de programació molt utilitzat en l’àmbit de la ciència de dades i que té l’avantatge de ser entenedor i senzill. Gràcies a biblioteques com \texttt{requests} (per descarregar dades d’una API) o \texttt{json} (per treballar amb el format de les dades), vaig poder construir el codi necessari.
% sense complicacions excessives.

Un altre avantatge de Python és que el seu codi es pot llegir amb facilitat, fins i tot per persones que no són expertes. Això facilita molt l’explicació del treball.

\section{Mobilitat sostenible}
La part final de la recerca prèvia va consistir en contextualitzar la importància d’aquest tema. El transport és una de les principals fonts de contaminació a les ciutats, i serveis com el Bicing contribueixen a reduir emissions i fomentar hàbits saludables.

Barcelona, a més, és una ciutat que aposta fort per la bicicleta: cada any s’amplien els carrils bici i es potencien mesures per animar la ciutadania a pedalar. Això fa que aquest treball no sigui només una experiència tècnica, sinó també una contribució dins un repte social de gran actualitat: la mobilitat sostenible.

\section{Iniciatives de dades obertes en mobilitat}
Més enllà de l’API oficial del Bicing, hi ha projectes i iniciatives que han tingut un paper important
en la difusió de les dades obertes aplicades a la mobilitat. Un dels casos més destacats és
\textbf{CityBikes}~\cite{citybikes}, una plataforma creada per un desenvolupador independent
que agrega informació de serveis de bicicleta compartida d’arreu del món.
El projecte posa a disposició una API pròpia que centralitza dades de diferents ciutats,
fent més fàcil construir aplicacions globals o comparar serveis de bicicletes entre països.

A escala local, la \textbf{Fundació Iniciativa Barcelona Open Data}~\cite{iniciativa-opendata}
és una entitat que promou la divulgació i reutilització de dades obertes a Barcelona i la seva
àrea metropolitana. Aquesta fundació treballa per sensibilitzar la ciutadania, oferir formació
i impulsar l’ús responsable de la gran quantitat d’informació que les administracions i serveis
fan pública.

Aquestes iniciatives mostren que les dades obertes no només són una eina tècnica,
sinó també una oportunitat social i econòmica. Ara bé, cal remarcar que la disponibilitat de dades
comporta la responsabilitat d’interpretar-les correctament i utilitzar-les de manera ètica.

